\chapter{Prefacio: por qué la economía necesita números}
\label{cap:prefacio}

\section{La gestión es decidir con información}

Gestionar una organización —una empresa, una cooperativa, un hospital, un Estado— es, en el fondo, \textbf{tomar decisiones bajo escasez}. Hay un presupuesto limitado, una capacidad de producción limitada, un tiempo limitado y una demanda que reacciona a lo que hacemos. La economía es precisamente la disciplina que estudia cómo se asignan recursos escasos entre fines alternativos, y por eso ofrece un marco natural para pensar la gestión.

Ahora bien, durante mucho tiempo esas decisiones se tomaron «a ojo», apoyadas en la experiencia. El cambio de las últimas décadas es que hoy disponemos de \textbf{datos} y de \textbf{herramientas de cálculo} que permiten cuantificar lo que antes era intuición. Saber que «si subo el precio, vendo menos» es economía básica; saber \emph{cuánto} menos vendo, a partir de los datos de mi propio negocio, es economía cuantitativa aplicada. Esa diferencia —pasar del signo a la magnitud— es el corazón de este curso.

\begin{ideabox}
\textbf{Idea central del Laboratorio.} La teoría económica nos dice \emph{qué} relación esperar entre variables (precio y cantidad, costo y producción, inversión y rentabilidad); las herramientas computacionales nos dicen \emph{cuánto} vale esa relación con los datos concretos que tenemos enfrente. Una sin la otra queda coja.
\end{ideabox}

\section{Qué vas a encontrar en este material}

Este texto recorre, en orden, los bloques económicos del Laboratorio:

\begin{enumerate}
  \item \textbf{Datos, variables y visualización} (Capítulo~\ref{cap:datos}): el punto de partida práctico.
  \item \textbf{Funciones económicas} (Capítulo~\ref{cap:funciones}): demanda, oferta, costos, ingreso y beneficio.
  \item \textbf{Equilibrio de mercado y sistemas de ecuaciones} (Capítulo~\ref{cap:equilibrio}).
  \item \textbf{El modelo insumo-producto de Leontief} (Capítulo~\ref{cap:leontief}).
  \item \textbf{Análisis marginal y derivadas} (Capítulo~\ref{cap:marginal}).
  \item \textbf{Elasticidades} (Capítulo~\ref{cap:elasticidad}).
  \item \textbf{Integrales y excedentes} (Capítulo~\ref{cap:integrales}).
  \item \textbf{Optimización económica} (Capítulo~\ref{cap:optimizacion}).
  \item \textbf{Competencia imperfecta: el duopolio} (Capítulo~\ref{cap:duopolio}).
  \item \textbf{Evaluación de proyectos de inversión} (Capítulo~\ref{cap:inversiones}).
\end{enumerate}

Cada capítulo sigue la misma estructura pensada para quien recién arranca: primero la \emph{intuición} (qué problema de gestión resolvemos), después la \emph{definición formal} y las fórmulas, luego un \emph{ejemplo numérico} resuelto a mano, y al final una caja \textsf{Del concepto al código} que muestra cómo ese mismo concepto reaparece —ya como herramienta— en las \emph{notebooks} de Python del laboratorio.

\section{Una aclaración sobre los prerrequisitos}

Esta materia combina dos mundos: la \textbf{economía} y la \textbf{programación}. Es posible que vengas con base en uno y no en el otro, o en ninguno. Por eso el material parte de cero en lo económico: no damos por sabido qué es una función de demanda ni qué significa «maximizar el beneficio». Si algún concepto te resulta nuevo, es esperable; está explicado dentro del propio texto y referenciado a manuales clásicos por si querés profundizar.

Lo que sí conviene tener fresco son tres herramientas matemáticas que usaremos como lenguaje: las \emph{funciones} (cómo una variable depende de otra), las \emph{derivadas} (cómo medir un cambio) y las \emph{integrales} (cómo acumular). Las tres se repasan en los capítulos donde hacen falta, siguiendo el enfoque de Chiang y Wainwright~\cite{chiang} y de Sydsæter y Hammond~\cite{sydsaeter}, pensados justamente para economistas que no son matemáticos profesionales.

\section{Cómo leer las fórmulas sin asustarse}

Una fórmula económica es una frase escrita en taquigrafía. Cuando veas, por ejemplo,
\[
\Pi(q) = I(q) - C(q),
\]
leelo en castellano: «el beneficio $\Pi$ que obtengo al producir una cantidad $q$ es igual al ingreso $I$ que genera esa cantidad, menos el costo $C$ de producirla». Toda la matemática del curso es de este tipo: traducciones precisas de ideas económicas razonables. Si entendés la idea, la fórmula deja de dar miedo.

Con esa actitud —entender antes que memorizar— arrancamos.
